\begin{abstract}
\noindent Software companies are often judged by the Rule of 40: revenue growth plus profit margin should reach 40\%. We show that the rule draws a line of equal score, not of equal value, and ask how high the score can go, even if AI does all the work. When capacity is built ahead of revenue, a company's best score is the margin it earns growing with its market, plus that market's growth net of the capital growth needs, plus a gain from taking share that cannot last. Under a verification hypothesis, a value-maximizing company weighs a point of growth at the value of an extra dollar of revenue, above the two to three points of margin investors' screens use when its revenue multiple is well above three. In SEC filings as first reported by \VFirmsEone{} software companies, 2010--25, the score that marks the top 5\% fell from 2016--19 to 2023--25 by \TopFallOpEone{} points on GAAP operating margins (\ScTopOpCoreFallLo{} to \ScTopOpCoreFallHi{} across samples; significant at 10\% under \TopFamTen{} of four corrections for multiple tests) and by \TopFallCashEone{} on free cash flow after stock-based pay (not significant), as fewer companies grew fast. Among top companies, a point more growth came with about a fifth of a point less margin; the data cannot show the weight companies act on. At central costs and fixed prices, full automation keeps a mature, asset-light company's cash score below 100 unless its market grows over \GstarCenOne\% a year, with AI compute at \CcCen\% of revenue (\GstarCcLo{} to \GstarCcHi\% across compute costs; \GstarPhiHalf{} or \GstarPhiSeven\% if half or 70\% of savings pass to customers). In a central scenario, not a forecast, a top company realizes \PaceCenSix\% of its program's benefits by 2036 and reaches a lasting score of \CenRSix{} (\CenRSixNet{} while still paying for changes under way; \CenRSixLo{} to \CenRSixHi{} from sampling alone), and repays the program's costs in \PayBackCen{} years. Acquisitions and interest distort reported scores. With free entry, automation lifts only companies with an input others cannot copy, up to a ceiling set by AI's full costs.
\end{abstract}
